\section{Проверка}
\label{sec:verification}
\porozhdeno{\Claim{V01}{Проверка исходного сравнения отделена от его расчёта: она сверяет частоты и назначения с артефактами QIIME, полноту идентификаторов, роли образцов, суммы категорий и расстояния. Испытания используют малый пример с ручными ответами и искусственно повреждённые выходы; локальное воспроизведение сравнивается с сохранёнными таблицами. Для семейственных профилей независимо проверяются маски отбора, суммы и отображаемые группы, включая отказ при подмене рида.}}

\porozhdeno{\Claim{V02}{Точное совпадение трёх технических исполнений DADA2 подтверждает повторяемость вычислений на фиксированном входе. Однако полный верификатор денойзинга сохраняет \texttt{INCOMPLETE/2}: у Kr5\_S84 осталось \ProjectKrFiveRetained{} из \ProjectKrFiveInput{} ридов (\ProjectKrFiveRetentionPct{}\%), ниже установленного критерия сохранности. Отдельный \texttt{PASS} проверки точного совпадения с согласованным исключением Kr5 этого ограничения не отменяет.}}

\porozhdeno{\Claim{V03}{Перенос результатов в числовой комплект проверен независимо от генератора: выполнены \ArticleTransferChecks{} сверок источников, знаменателей, округления и значений TeX. В \ArticleTransferCases{} испытаниях командной строки обнаружены все \ArticleDetectedMutations{} содержательных подмен; исходный и восстановленный комплекты приняты. Эти проверки устанавливают соответствие сохранённым данным, а не истинную таксономию; смысл утверждений статьи дополнительно сопоставлен с первичными источниками и картой оснований.}}
